\subsection{符号与词}

* 设 S 是一个非空集合，其元素将被称为符号（这一术语适合我们心目中的元数学应用）。设 \(\mathrm{L}_{0}(\mathrm{S})\) 是由 S 生成的自由半群； \(\mathrm{L}_{0}(\mathrm{S})\) 的元素被称为词，并与 S 的元素组成的有限序列 \(A = (s_{i})_{0 \leqslant i \leqslant n}\) 等同。\(\mathrm{L}_{0}(\mathrm{S})\) 中的合成法则将采用乘法形式书写，因此AB是通过将A和B并置而得到的序列。空词 \(\phi\) 是 \(\mathrm{L}_{0}(\mathrm{S})\) 的恒等元。我们回顾一下，长度 \(l(A)\) （词 \(A \in \mathrm{L}_{0}(\mathrm{S})\) 的长度）是序列A中元素的数量；因此 \(l(AB) = l(A) + l(B)\) ，长度为1的词是符号。设 \(\mathrm{L}(\mathrm{S})\) 表示 \(\mathrm{L}_{0}(\mathrm{S})\) 中的非空词所成的集合 .

此外，假设我们还给定了一个映射 \(s \to n(s)\) ，它把 \(S\) 映入集合 \(\mathbf{N}\) ——即整数 \(\geqslant 0\) 所组成的集合——。对于每个非空词 \(A = (s_i)_{0 \leqslant i \leqslant k}\) （ \(\mathbf{L}(\mathbf{S})\) 中的），我们置

\[
n (A) = \sum_ {i = 0} ^ {k} n (s _ {i})
\]

并且 \(n(\phi) = 0\) . \(n(A)\) 被称为 \(A\) 的权重。显然 \(n(AB) = n(A) + n(B)\) 。如果 \(A = A'BA''\) ，则词 \(B\) 被称为 \(A\) 的一个段（若还有 \(B \neq A\) ，则称为真段）。如果 \(A'\) （相应地， \(A''\) ）为空， \(B\) 被称为 \(A\) 的初始段（相应地，末段）。如果 \(l(A') = k\) ，那么 \(B\) 被称为始于第 \((k + 1)\) 位。

如果 \(A = BCDEF\) （其中词 \(B, C, D, E, F\) 可能为空），则段 \(C\) 和 \(E\) 被称为 \(A\) 的不相交的段。

